
\subsection*{7.16. Passage aux fibres géométriques.}
Soit @@M0@@ une clôture algébrique de @@M1@@, limite inductive de ses sous-extensions finies @@M2@@, et soit @@M3@@ le normalisé de @@M4@@ dans @@M5@@, limite inductive de ses normalisés @@M6@@ dans les @@M7@@. Choisissons un idéal maximal @@M8@@ de @@M9@@, induisant donc sur chaque @@M10@@ un idéal maximal @@M11@@, et soit @@M12@@, qui par Krull--Akizuki est encore un anneau de valuation discrète, dont nous désignons le spectre par @@M13@@, les points par @@M14@@ et @@M15@@. Le corps résiduel @@M16@@ est une clôture algébrique de @@M17@@, et c'est la limite inductive des corps résiduels @@M18@@ des @@M19@@. Soit @@M20@@, @@M21@@. \emph{Supposons maintenant} @@M22@@ \emph{de type fini sur} @@M23@@, donc @@M24@@ et @@M25@@ de type fini sur @@M26@@ resp. @@M27@@, donc \emph{noethériens}. En vertu de IV 3.2.3, on a
@@M28@@
D'autre part, en vertu de 7.15 on a un homomorphisme du système inductif des @@M29@@ dans le système inductif des @@M30@@. Passant à la limite, on trouve un homomorphisme canonique :
@@M31@@
appelé encore, comme de juste, \emph{homomorphisme de spécialisation}. Il induit un homomorphisme @@M32@@ grâce à 7.4, en utilisant le fait (dont la vérification est laissée au lecteur) que le foncteur @@M33@@ (tout comme @@M34@@) commute aux limites qui interviennent ici. De même, lorsque @@M35@@ est lisse, on peut interpréter (7.16.1) comme un homomorphisme sur les @@M36@@, qui sera d'ailleurs un homomorphisme d'anneaux, et même un homomorphisme de @@M37@@-anneaux si @@M38@@ est séparé~; dans ce cas, il induit donc un homomorphisme sur les @@M39@@.

